
% ------------------------------------------------------------
\section{Derivation of the Fokker--Planck equation}
\label{sec:derivation}
% ------------------------------------------------------------

\subsection{From a lattice random walk to a PDE}

Consider a symmetric random walk on a 1D lattice with spacing $h$ and a
position-dependent jump rate $\lambda(x)$.  Let $p(x,t)$ be the probability of
finding the particle at site $x$ at time $t$.  Starting from a left/right
symmetric rule,
\[
x \longrightarrow x\pm h \quad\text{at rate}\quad \lambda(x\pm h),
\]
the master equation is
\begin{equation}
\frac{\partial p(x,t)}{\partial t}
= \lambda(x+h)\,p(x+h,t) + \lambda(x-h)\,p(x-h,t) - 2\lambda(x)\,p(x,t).
\label{eq:master}
\end{equation}
Here the jump rate entering each incoming term is evaluated \emph{at the
departure site}; a particle leaves $x\pm h$ at rate $\lambda(x\pm h)$.
Taylor expanding $\lambda(x\pm h)p(x\pm h,t)$ about $x$,
\[
\lambda(x\pm h)\,p(x\pm h,t)
= \lambda(x)\,p(x,t) \pm h\,\partial_x[\lambda(x)\,p(x,t)]
+ \tfrac{h^2}{2}\,\partial_x^2[\lambda(x)\,p(x,t)] + O(h^3),
\]
so that the first-order terms cancel upon summation and\\[-.5ex]
\begin{equation}
\frac{\partial p}{\partial t}
= h^2\,\partial_x^2\!\bigl[\lambda(x)\,p(x,t)\bigr] + O(h^4).
\label{eq:nabla2lp}
\end{equation}
Taking the diffusive continuum limit $h\to 0$, $\lambda\to\infty$ with
$D(x)\;=\;\lambda(x)\,h^2$ held fixed yields the \emph{It\^o form} Fokker--Planck
equation,
\begin{equation}
\frac{\partial p}{\partial t} \;=\; \frac{\partial^2}{\partial x^2}\bigl[D(x)\,p\bigr].
\label{eq:ito-fpe}
\end{equation}

\subsection{Symmetric exclusion and Fick's form}

The It\^o form \eqref{eq:ito-fpe} follows from a rule in which the jump rate
depends on the \emph{departure} site.  In symmetric exclusion process (SEP), volume exclusion from the crowders restricts the \emph{arrival} site instead: a particle
at $x$ can hop to $x\pm h$ only if $x\pm h$ is empty.  If the probability that a
site is occupied by a crowder is $\rho(x)$, then the effective outgoing rate is
proportional to the fraction of empty neighbouring sites, $1-\rho(x\pm h)$,
times the intrinsic hop rate $\lambda_0$.  The master equation then takes the
equivalent ``gradient'' form\\[-.5ex]
\[
\frac{\partial p}{\partial t}
=\lambda_0\Bigl\{[1-\rho(x)]\,p(x+h,t) + [1-\rho(x)]\,p(x-h,t) - [1-\rho(x+h) + 1-\rho(x-h)]\,p(x,t)\Bigr\},
\]
whose continuum limit, after an analogous Taylor expansion, yields
\emph{Fick's form}:
\begin{equation}
{\;
\frac{\partial p}{\partial t}
= \frac{\partial}{\partial x}\!\left[D(x)\,\frac{\partial p}{\partial x}\right],
\quad D(x)=\lambda_0 h^2\,[1-\rho(x)].
\;}
\label{eq:fick-fpe}
\end{equation}
The two forms differ in the equilibrium distribution they produce in the absence
of flux.  The It\^o form has $p_\infty(x)\propto 1/D(x)$ (probability accumulates
in regions of low diffusivity); Fick's form has uniform
$p_\infty(x)=\mathrm{const}$, consistent with the SEP picture in which crowders
do not bias the \emph{density} of labelled particles.  Since our biological
motivation is foraging in a nest where a tracer bee should be uniform in the
absence of a food source, we adopt Fick's form throughout.



%Fick's form is chosen because it's the one where crowders slow you down without telling you where to go - which is exactly what should happen when there's no food source.
%The Itô form would give p_inf(x) \propto 1/D(x), meaning the bee accumulates where diffusivity is low (crowded regions).
%The Problem is: in the absence of food, the bee would appear to "prefer" crowded corridors - not because there's any food signal pulling it there, but purely as a mathematical artifact of how the noise is interpreted.
%This means when you later add a food source and try to detect where the bee spends more time, your signal is contaminated:
%
%Is the bee there because of food, or because of spurious drift from crowding?

A further practical advantage of Fick's form is that the operator
$\mathcal{L}p=\partial_x[D(x)\,\partial_x p]$ is self-adjoint on a periodic
domain: indeed, for any two smooth functions $f,g$ on $[-\pi,\pi]$ with periodic
data,
\[
\int_{-\pi}^{\pi} f\,\mathcal{L}g\;dx
=\int_{-\pi}^{\pi} f\,(Dg')'\;dx
=\underbrace{[f\,Dg']_{-\pi}^{\pi}}_{\substack{=0\\\text{periodic}}}
-\int_{-\pi}^{\pi} f'\,D\,g'\;dx,
\]
which is symmetric under $(f,g)\leftrightarrow(g,f)$.  Hence
$\mathcal{L}^\dagger=\mathcal{L}$ on periodic $L^2(-\pi,\pi)$.  For the It\^o
form the adjoint would differ by an extra $D'(x)T'(x)$ term, which is why the
It\^o SDE carries a spurious noise-induced drift (see the next subsection).

\subsection{The It\^o SDE associated with Fick's FPE}
\label{sec:sde}

The standard Fokker--Planck equation in drift--diffusion form is
\[
\partial_t p = -\partial_x(\mu p) + \tfrac12\,\partial_x^2(\sigma^2 p),
\]
corresponding to the It\^o SDE $dx = \mu(x)\,dt + \sigma(x)\,dW_t$.  Expanding
\[
\partial_x^2(\sigma^2 p)=(\sigma^2)''\,p+2(\sigma^2)'p'+\sigma^2 p'',\\[-.5ex]
\]
and matching coefficients with Fick's form $\partial_t p=Dp''+D'p'$ gives:
\begin{itemize}[noitemsep]
\item coefficient of $p''$: $\sigma^2/2 = D$, so $\sigma=\sqrt{2D}$;
\item coefficient of $p'$: $-\mu + (\sigma^2)'= -\mu+2D'=D'$, so $\mu=D'$;
\item coefficient of $p$: $-\mu'+\tfrac12(\sigma^2)''=-D''+D''=0$.
\end{itemize}
The It\^o SDE corresponding to Fick's form is therefore\\[-.5ex]
\begin{equation}
{\;
dx_t \;=\; D'(x_t)\,dt \;+\; \sqrt{2\,D(x_t)}\,dW_t.
\;}\label{eq:ito-sde}
\end{equation}
The noise-induced drift $D'(x)$ is a consequence of the spatial dependence of
$D$: without it, a naive discretisation of $dx=\sqrt{2D(x)}dW$ would produce the
wrong stationary distribution.  For the tent profile \eqref{eq:tent},
$D'(x) = -(D_1/\pi)\,\mathrm{sign}(x)$ on $[-\pi,\pi]\setminus\{0\}$ with a
jump at $x=0$.

\subsection{Domain and boundary conditions}
\label{sec:bcs}

We work on the periodic ring $x\in[-\pi,\pi]$ (so $x=-\pi$ is identified with
$x=\pi$), with a single partially absorbing trap at $x=0$.  The trap is modelled
by a $\delta$-sink of strength $\kappa>0$, so the full forward FPE reads
\begin{equation}
\partial_t p
=\partial_x\!\bigl[D(x)\,\partial_x p\bigr]
-\kappa\,\delta(x)\,p,
\qquad
p(-\pi,t)=p(\pi,t),\quad
\partial_x p(-\pi,t)=\partial_x p(\pi,t).
\label{eq:forward}
\end{equation}
The parameter $\kappa$ has units of velocity; $1/\kappa$ is the characteristic
dwell time at the trap before absorption.  The limit $\kappa\to\infty$ recovers
a fully absorbing trap (Dirichlet condition at $x=0$), while $\kappa\to 0$ gives
the trap-free ring (total probability conserved).
